% Part: normal-modal-logic
% Chapter: syntax-and-semantics
% Section: introduction

\documentclass[../../../include/open-logic-section]{subfiles}

\begin{document}

\olfileid{nml}{syn}{int}

\olsection{పరిచయం}

మోడల్ తర్కం \emph{మోడల్ ప్రతిజ్ఞావాక్యాలను},
వాటి మధ్య అనుగమన సంబంధాలను పరిశీలిస్తుంది.
అటువంటి వాక్యాలకు కొన్ని ఉదాహరణలు:
\begin{enumerate}
\item $2+2=4$ అవడం తప్పనిసరి.
\item రేపు వర్షం పడే సాధ్యత ఉండడం తప్పనిసరి.
\item $!A$ సాధ్యమై ఉండడం తప్పనిసరి అయితే,
  $!A$ సాధ్యం.
\end{enumerate}
సాధ్యత, అవశ్యకత మాత్రమే మోడాలిటీలు కావు:
ఇతర ఏకస్థానిక సంయోజకాలనూ మోడాలిటీలుగా
వర్గీకరిస్తారు. ఉదాహరణకు, “$!A$ జరగాలి,”
“$!A$ జరుగుతుంది,” “$!A$ అని డానాకు తెలుసు,”
లేదా “$!A$ అని డానా నమ్ముతుంది.”

అరిస్టాటిల్ రాసిన \emph{De Interpretatione}లో
మోడల్ తర్కం మొదట కనిపిస్తుంది. అవశ్యకమైనది
సాధ్యమూ అవుతుందని, కానీ సాధ్యమైన ప్రతిదీ
అవశ్యం కాదని; సాధ్యతను, అవశ్యకతను ఒకదానితో
మరొకటి నిర్వచించవచ్చని; $!A \land !B$
సాధ్యంగా సత్యమైతే $!A$, $!B$ విడివిడిగా
సాధ్యంగా సత్యమవుతాయని, కానీ తిరుగు దిశ
నిలవదని; $!A \to !B$ అవశ్యమైతే, $!A$
అవశ్యమైనప్పుడు $!B$ కూడా అవశ్యమని
మొదట గమనించింది ఆయనే.

ఆధునిక మోడల్ తర్కానికి తొలి విధానం
సి.~ఐ. లూయిస్ పనిలో ఏర్పడి,
లూయిస్, లాంగ్‌ఫర్డ్‌ల \emph{Symbolic Logic}
(1932)తో పరిపూర్ణమైంది. నిగమన సంబంధాన్ని
ద్రవ్యపర సోపాధికంతో సూచించడం వారికి
సంతృప్తికరంగా అనిపించలేదు:
“$!A$ నుంచి $!B$ అనుగమిస్తుంది”
అనడానికి $!A \lif !B$ సరైన ప్రత్యామ్నాయం
కాదు. బదులుగా “$!A$ అయితే $!B$
అనేది అవశ్యం” అని నిగమన సంబంధాన్ని
వర్ణించాలని వారు ప్రతిపాదించారు;
దానికి $!A \strictif !B$ సంకేతం ఇచ్చారు.
ఈ లక్షణాలను వేరు చేస్తూ లూయిస్ ఐదు
మోడల్ వ్యవస్థలను, \Log{S1}, \ldots,
\Log{S4}, \Log{S5}, గుర్తించాడు;
చివరి రెండు నేటికీ వాడుకలో ఉన్నాయి.

లూయిస్, లాంగ్‌ఫర్డ్‌ల విధానం పూర్తిగా
\emph{వాక్యనిర్మాణపరమైనది}: వారు
సమంజసమైన స్వీకృతాలు, నియమాలను
గుర్తించి, వాటితో ఏమి నిరూపించవచ్చో
పరిశీలించారు. అర్థపర విధానం చాలాకాలం
అందుబాటులోకి రాలేదు. రుడోల్ఫ్ కార్నాప్
\emph{Meaning and Necessity} (1947)లో
\emph{స్థితి వివరణ} అనే భావనతో తొలి
ప్రయత్నం చేశాడు. అది ఆ స్థితి
వివరణలో “సత్యం” అయ్యే పరమాణు
వాక్యాల సమాహారం. ఏ వాక్యం $!A$కైనా
సత్య నిర్వచనాన్ని విస్తరించిన తరువాత,
అన్ని స్థితి వివరణల్లో సత్యమయ్యే
$!A$ను కార్నాప్ \emph{అవశ్యంగా సత్యం}
అని నిర్వచించాడు. అయితే
“సాధ్యంగా అవశ్యంగా \ldots{}
సాధ్యంగా $!A$” వంటి వాక్యాలు
ఎప్పుడూ అంతర్గత మోడాలిటీకి తగ్గిపోవడంతో,
కార్నాప్ విధానం \emph{పునరావృత}
మోడాలిటీలను నిర్వహించలేకపోయింది.

మోడల్ అర్థవిచారంలో కీలక పురోగతి
సాల్ క్రిప్కె “A Completeness Theorem in
Modal Logic” (JSL 1959) వ్యాసంతో వచ్చింది.
ఒక ప్రకటన “సాధ్యమైన అన్ని లోకాల్లో”
సత్యమైతే అవశ్యంగా సత్యమనే లైబ్నిజ్
ఆలోచనపై క్రిప్కె తన పనిని ఆధారపరిచాడు.
కానీ ఒక లోకం $w$లో (లేదా స్థితి
వివరణ $s$లో) ప్రకటన సత్యం కావడం
$w$పై అసలు ఆధారపడకపోవడం వల్ల,
ఈ ఆలోచనకూ కార్నాప్ విధానంలోని
లోపమే ఉంది. కాబట్టి లోకాలు
\emph{ప్రాప్యత సంబంధం} $R$తో
కలిసి ఉన్నాయని క్రిప్కె అనుకున్నాడు.
“అవశ్యంగా $!A$” అనే ప్రకటన
లోకం $w$లో సత్యం అవుతుంది,
$w$ నుంచి \emph{ప్రాప్యమైన}
అన్ని లోకాలు $w'$లో $!A$
సత్యమైతే, అప్పుడు మాత్రమే.
ఈ విధానపు ఏదో రూపాన్ని ఇచ్చే
అర్థవిచారాలను క్రిప్కె అర్థవిచారాలు
అంటారు; అవి మోడల్ తర్కాల విస్తృత
అభివృద్ధికి మార్గం తెరిచాయి.

క్రిప్కె అర్థవిచారంలో చదివినప్పుడు,
\emph{సంబంధ నిర్మాణాలు} “లోపలి నుంచి”
ఎలా కనిపిస్తాయో మోడల్ తర్కం చూపిస్తుంది.
సంబంధ నిర్మాణం అంటే ద్విస్థానిక
సంబంధంతో కూడిన సమితి మాత్రమే;
ఉదాహరణకు, సామాజిక భద్రతా సంఖ్యల
క్రమంలో ఉంచిన ఒక తరగతి విద్యార్థుల
సమితి అలాంటి నిర్మాణం. అయితే
సంబంధ నిర్మాణాల పరిధులు అనేకం:
ప్రపంచ స్థితుల సాపేక్ష సాధ్యతతో
పాటు, ఒక వ్యక్తి జ్ఞాన సంబంధిత
స్థితులు జ్ఞానపర సాధ్యత ద్వారా
సంబంధించి ఉండవచ్చు; చలన వ్యవస్థ
స్థితులు వాటి మార్పుల ద్వారా
సంబంధించి ఉండవచ్చు. వీటన్నిటినీ
మోడల్ తర్కంతో నమూనా చేయవచ్చు:
మొదటిది సాధారణ అలెథిక్ (సత్య-అవశ్యకత)
మోడల్ తర్కాన్ని, మిగతావి జ్ఞానమీమాంసాత్మక
తర్కం, గతిక తర్కం మొదలైనవాటిని ఇస్తాయి.

మోడల్ తార్కికులకు “అనురూపతా సిద్ధాంతం”గా
తెలిసిన ఒక ప్రత్యేక కోణంపై ఇక్కడ దృష్టి
పెడతాం. క్రిప్కె తొలి ముఖ్య
ఆవిష్కరణల్లో ఒకటి: ప్రాప్యత
సంబంధం~$R$కి ఉన్న సంక్రమణశీలత,
సౌష్టవం వంటి అనేక లక్షణాలను
తగిన “మోడల్ పథకాల” ద్వారా
\emph{మోడల్ భాషలోనే} లక్షణీకరించవచ్చు.
ఉదాహరణకు, $R$ స్వప్రావర్తకత
“$!A$ అవశ్యమైతే $!A$”
అనే పథకానికి “అనురూపం”
అని మోడల్ తార్కికులు అంటారు.
\Ax{D}, \Ax{T}, \Ax{B}, \Ax{4},
\Ax{5} పథకాలను కలిపి పొందే
సాంప్రదాయిక మోడల్ తర్క వ్యవస్థల
(ఉదా., \Log{S4}, \Log{S5})
అనురూపతా సిద్ధాంతాన్నే
ముఖ్యంగా పరిశీలిస్తాం.

\end{document}
